\documentclass[11pt]{article}
\usepackage{../../../tex/lecturestyle}

\title{Wavelet Matrix (Basics)}
\author{}
\date{}

\begin{document}
\thispagestyle{fancy}

\lecturetitleband{Data Structures}{Wavelet Matrix (Basics)}{6}

\section{Overview}
Fix a static array $A = (A_0, \ldots, A_{N-1})$ of integers in
$[0, 2^H)$. A \emph{wavelet matrix} answers, for any query range
$[L, R)$, questions about the multiset $\{A_L, \ldots, A_{R-1}\}$ that a
binary trie built just from those elements could answer --- in
$O(H)$ time, without ever building that trie. This lecture covers two such
questions:

\begin{itemize}[leftmargin=1.5em]
  \item the $k$-th smallest value among $A_L, \ldots, A_{R-1}$ (0-indexed);
  \item the number of elements of $A_L, \ldots, A_{R-1}$ lying in $[a, b)$.
\end{itemize}

The construction generalizes well past these two queries once its
mechanics are clear; a follow-up lecture on further applications is out of
scope here.

\section{Binary Tries, Briefly}
A height-$H$ binary trie over a multiset of integers in $[0, 2^H)$ is the
binary tree obtained by classifying each integer by its bits, most
significant first. A node at height $h$ (the root has height $H$, leaves
height $0$) corresponds to an interval $[l, r)$ with $r - l = 2^h$; its left
child holds the elements whose next bit (the $(h-1)$-th) is $0$, its right
child those whose bit is $1$.

\begin{definitionbox}
If every node stores only the \emph{count} of elements below it, two queries
become easy in $O(H)$ time by walking root to leaf:
\begin{itemize}[leftmargin=1.5em]
  \item \textbf{$k$-th smallest} (0-indexed): at each internal node, let $x$
        be the left child's count. If $k < x$, descend left; otherwise
        descend right with $k \leftarrow k - x$.
  \item \textbf{count $< a$}: track the node's integer interval $[l, r)$ as
        you descend, with $m = (l+r)/2$ splitting it. If $a \le m$, all
        qualifying elements are in the left child, so descend left.
        Otherwise every element in the left child is $< a$: add the left
        child's count to a running answer and descend right. A range count
        on $[a, b)$ is then (count $< b$) $-$ (count $< a$).
\end{itemize}
\end{definitionbox}

Building one trie for each of $O(N)$ possible query ranges is obviously too
slow. The wavelet matrix's contribution is to let \emph{one} structure,
built once in $O(NH)$, answer the trie query above for \emph{any} range
$[L, R)$ on demand.

\section{From Wavelet Tree to Wavelet Matrix}

\subsection{The key stability observation}
Suppose each trie node stores its elements \emph{in their original index
order} rather than unordered. Then restricting the whole array to a range
$[L, R)$ restricts every node's stored elements to some contiguous
sub-range of what it would otherwise hold --- because a node's content is a
stable sub-selection of its parent's content (elements keep their relative
order as they're partitioned by bit), a contiguous range of indices at the
root stays contiguous after any prefix of bit-tests. This is the whole idea
behind a \emph{wavelet tree}: a binary trie augmented, at every internal
node, with enough information (a prefix-sum of which child each element
went to) to compute that node's sub-range for any $[L, R)$ given the
parent's sub-range.

\subsection{Flattening it}
A wavelet tree stores this per-node information node by node. A
\emph{wavelet matrix} rearranges the same information level by level
instead: at height $h$, concatenate every node's ``did this element go
left or right'' bit sequence, in left-to-right node order, into one array
of length $N$. Because sibling nodes' element blocks are contiguous and
appear in a fixed left-then-right order at each level, this reshuffling
loses nothing --- the position of the interval boundaries between nodes at
a given height is itself recoverable from the data, exactly the way a node's
integer interval in a plain trie is recoverable from the path taken to
reach it.

The payoff is a much simpler physical layout: $H$ bit-arrays of length $N$,
each with a rank (prefix-sum) table, and nothing else.

\subsection{The core operation}
Represent a trie node by a triple $(h, [l, r))$: it is the node at height
$h$, and the elements currently assigned to it occupy positions $[l, r)$ in
the height-$h$ bit-array (equivalently, of the array as reordered after
$H - h$ levels of partitioning). Let $\mathrm{dat}_h[i]$ be the number of
$1$-bits among the first $i$ entries of the height-$h$ bit-array (its rank
table). Then the two children of $(h+1, [l, r))$ are computable in $O(1)$:
letting
\[
  a_1 = \mathrm{dat}_h[l],\quad a_0 = l - a_1, \qquad
  b_1 = \mathrm{dat}_h[r],\quad b_0 = r - b_1, \qquad
  c_0 = N - \mathrm{dat}_h[N],
\]
the left child is $(h, [a_0, b_0))$ and the right child is
$(h, [c_0 + a_1,\, c_0 + b_1))$.

\begin{notebox}
Why this formula: $a_0$ (resp.\ $b_0$) counts how many of the first $l$
(resp.\ $r$) elements at this level went left, so $[a_0, b_0)$ is exactly
where our slice's ``went-left'' elements land among all the went-left
elements globally --- and all went-left elements are stored first at height
$h$ (positions $[0, c_0)$, with $c_0$ the total left-count). The
went-right elements come right after, at $[c_0, N)$, so the same
counting argument applied to $a_1, b_1$ (now counting within the
went-right elements) gives the offset $c_0 + a_1$.
\end{notebox}

\section{Build and Query Algorithms}

\subsection{Build}
Process heights $h = H-1$ down to $0$. At each height, with the array in
its current order, read off the bit at position $h$ of every element,
stably partition the array into ``bit $0$'' elements followed by ``bit
$1$'' elements, and record the bit sequence's prefix sums as
$\mathrm{dat}_h$ before moving on. After all $H$ levels this is exactly a
most-significant-bit-first radix sort, and the total work is $O(NH)$.

\subsection{$k$-th smallest}
Starting from the root $(H, [L, R))$, at each step compute the current
node's two children via the formula above; if $k$ is less than the left
child's size, descend left, otherwise subtract the left size from $k$ and
descend right (recording a $1$ bit at this height in the answer). After $H$
steps the accumulated bits are the answer, found in $O(H)$ time.

\subsection{Range count in $[a, b)$}
Define $\mathrm{countLess}(l, r, x)$, the number of elements of the current
slice below $x$: walk from the root, and at each height compare $x$'s bit
to the branch just as in the plain-trie version --- if $x$'s bit is $1$,
every element that went left at this level is $< x$ regardless of lower
bits, so add the left child's size to the running total and continue into
the right child (comparing the remaining bits); if $x$'s bit is $0$,
nothing is added and the walk continues into the left child only. The
range-count answer is $\mathrm{countLess}(L, R, b) - \mathrm{countLess}(L,
R, a)$, again $O(H)$ total.

\begin{theorembox}
Both queries above run in $O(H)$ time after $O(NH)$-time, $O(NH)$-space
preprocessing.
\end{theorembox}

\section{A Note on rank}
Both algorithms lean on $\mathrm{dat}_h[i]$: a rank query, ``how many $1$s
in the first $i$ bits.'' The examples below use a plain prefix-sum array
per level, which is simple and already achieves the stated complexity.
Competitive-programming implementations more often pack each level's bits
into 64-bit words, store one running count per word (rather than per bit),
and answer a rank query by adding a hardware \texttt{popcount} of the
remaining partial word --- the same asymptotic $O(1)$ per query, with a
much smaller constant and a fraction of the memory. Taken further, with a
genuinely succinct rank/select structure on each bit-array, the whole
wavelet matrix can be stored in close to the information-theoretic minimum
$NH(1+o(1))$ bits; that refinement is rarely the bottleneck in contest
settings and is not covered here.

\section{Reference Implementations}

\subsection{Rust}
\lstinputlisting[language=Rust]{example.rs}

\subsection{Python}
\lstinputlisting[language=Python]{example.py}

\section{Worked Example}
\begin{examplebox}
Let $A = (3,1,2,1,1,0,3,2,2)$, so $H = 2$. For the range $[2, 7)$ (values
$2,1,1,0,3$), the sorted order is $0,1,1,2,3$, so the $0$-th smallest is
$0$ and the $4$-th smallest is $3$ --- matching the reference
implementations, which also cross-check every $(l, r, k)$ and $(l, r, a,
b)$ combination against a brute-force scan over the whole array.
\end{examplebox}

\section{Further Reading}
\begin{itemize}[leftmargin=1.5em]
  \item AtCoder Algorithm Lectures index (source of this curriculum):
        \url{https://info.atcoder.jp/entry/algorithm_lectures/index}
  \item F. Claude, G. Navarro, A. Ordóñez Pereira, ``The wavelet matrix: An
        efficient wavelet tree for large alphabets'', \emph{Information
        Systems} 47, 2015, pp.\ 15--32.
  \item USACO Guide, \emph{Wavelet Tree}: \url{https://usaco.guide/adv/wavelet}
\end{itemize}

\end{document}
